\section{The Difference Between a Difficult Problem and a Valuable Problem}

\subsection{The 3 Categories of Difficulty}

People are often told that the way to become wealthy is to solve difficult problems. This contains an important truth, but it is incomplete in a way that can make it dangerous for technically ambitious people.

Consider 3 problems.

One is to prove an obscure mathematical proposition. It may be extraordinarily difficult and intellectually beautiful, yet have no immediately identifiable buyer. Another is to help a hospital prevent a particular class of billing errors. The underlying software might be comparatively straightforward, but the implementation could be commercially valuable because the hospital loses money and time every month. The third is to develop a safe new drug. This is both technically difficult and potentially immensely valuable, but it requires research capital, regulatory approval, long time horizons, and the ability to survive failure.

The 3 problems differ along several dimensions that have almost nothing to do with the founder's intelligence: the size of the benefit, the existence of a buyer, the urgency of the need, the availability of alternatives, the cost of delivering a solution, the time to evidence, the amount of required capital, and the ability to appropriate some of the resulting value.

\textbf{Difficulty is not the objective. Economically consequential difficulty, coupled with a credible path to delivery and payment, is an opportunity.}

\subsection{Why Difficulty Protects but Does Not Guarantee Value}

Why, then, do difficult problems matter? Because they can create a gap between firms that can solve them and firms that cannot. If solving something requires uncommon scientific expertise, years of relationship-building, complex operational coordination, or a demanding combination of disciplines, competitors cannot necessarily copy the result by copying a website. Difficulty can protect an advantage once an actual customer benefit exists.

But there are at least 4 reasons why a difficult problem may still make a poor business:
\begin{enumerate}[label=\textbf{\arabic*.},leftmargin=11pt,itemsep=1pt,topsep=1.5pt,parsep=0pt]
  \item \textbf{The beneficiary and the buyer may be different.} An improved public-health method may help millions, but capturing payment for it may require government procurement, insurance reimbursement, or philanthropy.
  \item \textbf{Verification asymmetry.} The person who benefits may not know the benefit exists, or may be unable to verify it.
  \item \textbf{Excessive delivery cost.} A proposed solution may be too expensive to produce, even when the gross benefit is large.
  \item \textbf{Adoption friction.} Adoption may demand painful organizational changes that the customer is unwilling to undertake.
\end{enumerate}

A founder might create technology that saves an enterprise ₹10 crore annually in principle. Yet if using it requires 18 months of integration, exposes the company to security risk, and requires retraining hundreds of employees, the buyer's willingness to pay will not simply equal a convenient fraction of the advertised savings.

\textbf{The real problem is larger than the technical puzzle. It includes the customer's environment, incentives, fears, constraints, and decision process.}

\subsection{The 5 Tests of an Economically Attractive Problem}

We can describe an economically attractive problem with 5 questions:
\begin{itemize}[leftmargin=11pt,itemsep=1pt,topsep=1.5pt,parsep=0pt]
  \item \textbf{Consequence:} What measurable or intensely felt benefit arises when the problem is solved?
  \item \textbf{Demand:} Who cares enough to act, and who controls the money or resources needed to act?
  \item \textbf{Feasibility:} Can a solution be produced and delivered reliably, at acceptable cost and within a useful time frame?
  \item \textbf{Appropriability:} Can the organization receive enough of the benefit to finance its work, despite competition, imitation, and bargaining?
  \item \textbf{Persistence:} Will the problem and the organization's ability to solve it remain relevant long enough to justify the investment?
\end{itemize}

These are not 5 bureaucratic boxes to be checked before doing anything. They are 5 ways for a founder to avoid mistaking intellectual excitement for a viable opportunity.

Paul Graham's advice to build things that people want is valuable partly because it forces attention away from the founder's delight in a solution and toward the customer's experience of a problem. His separate advice to do things that do not scale reminds us that, early on, manual selling and intensive customer work can reveal what the real problem is.\footnote{Paul Graham, ``Do Things that Don't Scale'' (July 2013).} A founder may think they are building a scheduling tool and discover that the customer's actual difficulty is unreliable coordination among several departments. The useful product may be quite different from the original prototype.

A surprisingly useful question is: If I could solve this problem perfectly, who would be grateful enough to make an actual sacrifice in exchange for the solution? Money is one sacrifice, but time, attention, organizational change, and political capital are others. People may praise ideas for free. Their willingness to sacrifice reveals priorities.

\textbf{The best problem to attack is not necessarily the hardest one you can understand. It is a problem where your ability to make progress is unusually high relative to the value that progress would create, and where a sustainable business can be constructed around that progress.}
